% Ad Quintum Fratrem 1.4 --- headnote
% ad-quintum-fratrem-01-04, around the Nones of August 58 BC, written at Thessalonica

\textit{Cicero to his brother Quintus, written from
Thessalonica around the Nones of August 58 BC. Now that
Quintus is at Rome and exposed to the full reach of his
brother's enemies, the letter shifts from the lament of
Q. fr. 1.3 to the strategic reckoning. \textsection 1 is the
careful self-defence: Cicero asks Quintus not to read his
ruin as wickedness or crime, but as the failure of his too
cautious counsel --- since each of those nearest him,
closest, most familiar, ``either was afraid for himself or
envied me.''}

\textit{\textsection 3 names the new tribunician college that
will take office on 10 December: Sestius, Curius, Milo,
Fadius, Fabricius are or are hoped to be friendly. But
Clodius, even as a private citizen once his year ends, will
have the same crowd to call out at the public meetings; and
an \textit{intercessor} (a tribunician veto) will be
procured. \textsection 4 is Cicero's bitterest summing-up of
the spring of 58 --- the things that were \textit{not} put to
him as he set out: ``Pompey's sudden defection, the
alienation of the consuls, even of the praetors, the fear of
the tax-farmers, arms.'' He had been told instead that he
should be returning within three days at the head of the
highest glory. The letter closes with the bargain again ---
``I shall hold on to life so long as I think it of consequence
to you'' --- and the warning that, with Quintus exposed too,
the fight will be by lawsuits, not swords.}
